\section{\texorpdfstring{\(\mathfrak{S}\)-收敛的一致结构}{S-收敛的一致结构}}

记号。若 X 与 Y 是任意两个集，我们回忆 X 到 Y 的一切映射组成的集记作 \(\mathcal{F}(\mathbf{X};\mathbf{Y})\)，它可与乘积集 \(\mathbf{Y}^{\mathbf{x}}\) 等同（《集合论》，第 II 章，§ 5，第 2 号）。对 \(\mathcal{F}(\mathbf{X};\mathbf{Y})\) 的每个子集 H 与每个 \(x\in \mathbf{X}\)，我们将用 \(\mathrm{H}(x)\) 表示元素 \(u(x)\in \mathbf{Y}\)（当 \(u\) 取遍 H 时）组成的集。若 \(\Phi\) 是 \(\mathcal{F}(\mathbf{X};\mathbf{Y})\) 上的滤子基，我们用 \(\Phi (x)\) 表示由诸集 \(\mathrm{H}(x)\)（H 取遍 \(\Phi\)）组成的 Y 上的滤子基。最后，我们回忆，对每个 \(u\in \mathcal{F}(\mathbf{X};\mathbf{Y})\) 与 X 的每个子集 A，\(u|\mathrm{A}\) 表示 \(u\) 在 A 上的限制，它是 A 到 Y 中的映射；若 H 是 \(\mathcal{F}(\mathbf{X};\mathbf{Y})\) 的子集，则 H\textbar A 将表示限制 \(u|\mathrm{A}\)（函数 \(u\in \mathrm{H}\) 的）组成的集。

